\section{Motivation}
\label{sec:motivation}

Robotic grasping is a fundamental capability in automated manufacturing, logistics, and service robotics\cite{Todescato}. The ability to reliably grasp objects of
varying shapes, sizes, and surface properties remains a challenging problem, particularly in unstructured environments\cite{Dinakaran2022A}. Industrial applications such
as bin picking, assembly, and material handling require grippers that can operate robustly across a wide range of workpieces\cite{Kang2019Design}.

Among the commonly used end-effectors, suction-cup grippers and two-finger grippers represent two complementary approaches\cite{Courchesne2023A}. Suction grippers are
widely used for handling objects with smooth surfaces due to their simplicity and speed, whereas two-finger grippers provide flexibility in handling
objects with complex geometries. Designing optimal configurations for these grippers is non-trivial, as grasp performance depends on multiple
interacting factors such as geometry, contact distribution, force limits, and collision constraints\cite{Liu2024Topology}.

While the gripper configuration problems are well-motivated in themselves, the \textbf{primary objective of this thesis} is not simply to find the
best gripper configuration, but to \textbf{systematically apply and compare optimization algorithms} for solving these configuration problems. A large
number of optimization algorithms exist, spanning evolutionary strategies, decomposition-based methods, indicator-based selection, and Bayesian
optimization, yet their relative suitability for gripper design problems of this type is not well understood. Different algorithms may exhibit
substantially different convergence behaviour, solution quality, and robustness depending on the structure of the objective landscape and the
dimensionality of the search space. Identifying which algorithm class performs best for each gripper problem is therefore a practically relevant and
scientifically meaningful research question\cite{Doerterler2021Analyzing}.

To provide a rigorous and reproducible basis for this comparison, two gripper configuration problems are formulated as well-defined multi-objective
optimization tasks. In both cases, a diverse set of workpieces is incorporated directly into the problem formulation by treating each workpiece as a
separate objective. This design choice ensures that the optimized gripper configurations are not tailored to a single
geometry, but are instead robust across a range of workpiece shapes. It also increases the dimensionality of the objective space as the number of
workpieces grows, allowing the benchmarking study to probe algorithm behaviour under varying levels of problem complexity\cite{Guo2022Evolutionary}.

For the \textbf{multi-suction-cup gripper}, the objectives consist of the geometric stability of the suction base together with the grasp diversity
evaluated independently for each workpiece. These objectives are inherently conflicting: a larger suction base improves stability but increases the
risk of collisions with the workpiece, reducing the number of feasible grasps.

For the \textbf{two-finger gripper}, each workpiece contributes one objective, representing how well the optimized finger contour conforms to that
particular workpiece geometry. A finger contour that scores well across all workpieces simultaneously must generalize across diverse geometric
profiles, making the problem increasingly challenging as more workpieces are added\cite{Liu2024Topology}.

Together, these two problems provide a structured experimental basis for comparing optimization algorithms across varying objective space
dimensionalities and objective landscape characteristics, within the practically relevant domain of robotic gripper design.


% The \textbf{multi-suction-cup gripper} problem is cast as a bi-objective optimization task. The spatial arrangement of suction cups influences two
% competing performance aspects:

% \begin{enumerate} \item \textbf{Support Triangle Area:} The area of the geometric base formed by the suction cups determines grasp stability and
%     resistance to external disturbances. A larger base provides better force distribution and tilt resistance.

%     \item \textbf{Grasp Diversity:} The number, spatial location, and orientation of valid collision-free grasps achievable on a given workpiece. This
%     metric is inherently workpiece-dependent, as the gripper geometry must conform to the object's surface without collisions.
% \end{enumerate}

% These two objectives are inherently conflicting: a larger suction cup base improves stability but increases the likelihood of collisions with the
% workpiece, thereby reducing the number of feasible grasps. This conflict makes the problem a suitable benchmark for \textbf{multi-objective
% evolutionary algorithms} such as NSGA-II, NSGA-III, MOEA/D, SMS-EMOA, and CMOPSO, as well as \textbf{multi-objective Bayesian optimization} methods.

% The \textbf{two-finger gripper} problem is cast as a single-objective optimization task. The finger contour is parameterized by a spline defined
% through five points, three of which have variable vertical positions $h_{\mathrm{spline},i}$ ($i = 2, 3, 4$), along with the gripper nose length
% $l_{\mathrm{gn}}$ and depth $d_{\mathrm{gn}}$. The objective is to maximize the voxel-based proximity score $s_{\max}(\mathbf{q})$, computed via
% FFT-based convolution between the voxelized gripper and workpiece geometries. A configuration is valid only if the required finger opening
% corresponds to a physical closing motion. This five-dimensional, black-box, single-objective problem serves as a benchmark for
% \textbf{single-objective evolutionary algorithms} and \textbf{Bayesian optimization} methods.